\documentclass{beamer}
\usetheme{Madrid}
\usecolortheme{crane} 
\setbeamertemplate{navigation symbols}{}
\setbeamertemplate{caption}[numbered]
\usepackage{graphicx}      
\usepackage{amsmath,amssymb} 
\usepackage{booktabs}      
\usepackage{algorithm}     
\usepackage{algpseudocode} 
\usepackage{capt-of}       
\usepackage{hyperref}   
\logo{\includegraphics[height=1.3cm]{"/home/bodike/Desktop/B24CS013_Bodike Chaithali_10/iit-bhilai.jpg"}\hspace{0.5cm}}
\setbeamercolor{background canvas}{bg=blue!15}
\setbeamercolor{title}{bg=yellow, fg=black}
\setbeamercolor{frametitle}{bg=yellow, fg=black}
\setbeamercolor{footline}{bg=yellow, fg=black}
\setbeamercolor{block title}{bg=yellow!80, fg=black}  
\setbeamercolor{block body}{bg=yellow!20, fg=black} 
\setbeamercolor{block title alerted}{bg=red!70, fg=white}
\setbeamercolor{block body alerted}{bg=red!20, fg=black} 
\setbeamercolor{block title example}{bg=green!60, fg=white}
\setbeamercolor{block body example}{bg=green!20, fg=black}
\setbeamertemplate{blocks}[rounded][shadow=false]
\setbeamertemplate{footline}{%
	\leavevmode%
	\hbox{%
		\begin{beamercolorbox}[wd=.33\paperwidth,ht=3ex,dp=1ex,left]{footline}%
			\hspace{7mm} Anand, Arindam (IIT)
		\end{beamercolorbox}%
		\begin{beamercolorbox}[wd=.34\paperwidth,ht=3ex,dp=1ex,center]{footline}%
			Beamer Lab Activity - 3
		\end{beamercolorbox}%
		\begin{beamercolorbox}[wd=.33\paperwidth,ht=3ex,dp=1ex,right]{footline}%
			Beamer 2025 \hspace{5em} \insertframenumber/\inserttotalframenumber \hspace{2mm}
		\end{beamercolorbox}%
	}%
}

\title[Beamer Lab Activity - 3]{Software Tools and Technologies Laboratory\\Latex Lab Activity - 3: Beamer}
\author[Anand, Arindam]{Anand M. Baswade\inst{1} \and Arindam Chakraborty\inst{2}}
\institute[IIT]
{
	\inst{1}
	Faculty of Computer Science and Engineering\\Indian Institute of Technology, Bhilai
	\and
	\inst{2}
	M.Tech Student of Computer Science and Engineering\\Indian Institute of Technology, Bhilai
}
\date{\today}
\AtBeginSection[]{
	\begin{frame}<beamer>
		\frametitle{Table of Contents}
		\tableofcontents[currentsection]
	\end{frame}
}
\bibliographystyle{plain}
\begin{document}
	\frame{\titlepage}
	\begin{frame}
		\frametitle{Table of Contents}
		\tableofcontents
	\end{frame}
	\section{Introduction}
	\begin{frame}{Introduction to \LaTeX{} Beamer}
		\begin{itemize}
			\item \LaTeX{} Beamer is a powerful class for creating presentations.
			\item It allows:
			\begin{itemize}
				\item Professional-quality slides
				\item Mathematical equations
				\item Figures, tables, and references
				\item Animations and overlays
			\end{itemize}
			\item Advantages:
			\begin{itemize}
				\item Consistent formatting
				\item Easy handling of citations and references
				\item Cross-platform PDF output
			\end{itemize}
		\end{itemize}
	\end{frame}
	\section{Why Beamer?}
	\begin{frame}{Why Use Beamer?}
		\begin{itemize}[<+->]
			\item Beamer is designed for creating professional presentations directly in \LaTeX{}.
			\item Consistent and elegant design
			\item Full support for mathematical equations
			\item Easy management of references and citations
			\item Portable PDF output (works everywhere)
			\item Highly customizable (themes, colors, backgrounds, animations)
			\item Free and open-source — no licensing issues
		\end{itemize}
	\end{frame}
	\begin{frame}{Comparison with Other Tools}
		\begin{table}
			\centering
			\begin{tabular}{|p{0.3\textwidth}|p{0.3\textwidth}|p{0.3\textwidth}|}
				\hline
				\onslide<1->{Feature} & \onslide<1->{Beamer} & \onslide<1->{PowerPoint} \\
				\hline
				\onslide<1->{Equations} & \onslide<2->{Excellent} & \onslide<2->{Limited} \\
				\onslide<1->{Citations} & \onslide<3->{Automatic} & \onslide<3->{Manual} \\
				\onslide<1->{Portability} & \onslide<4->{PDF (stable)} & \onslide<4->{File-dependent} \\
				\onslide<1->{Cost} & \onslide<5->{Free} & \onslide<5->{Proprietary} \\
				\onslide<1->{Customization} & \onslide<6->{High (code-based)} & \onslide<6->{Moderate (GUI-based)} \\
				\hline
			\end{tabular}
		\end{table}
	\end{frame}
	\section{Features}
	\begin{frame}{Key Features of \LaTeX\ Beamer}
		\begin{block}{Presentation Design}
			\begin{itemize}
				\item<1-> Predefined professional themes for elegant slides
				\item<2-> Customizable colors, fonts, and backgrounds
				\item<3-> High-quality PDF output across platforms
			\end{itemize}
		\end{block}
		\begin{block}{Content Capabilities}
			\begin{itemize}
				\item<4-> Full support for complex mathematical equations
				\item<5-> Easy insertion of tables, figures, and graphics
				\item<6-> Step-by-step content reveal using overlays
			\end{itemize}
		\end{block}
		\begin{block}{References \& Automation}
			\begin{itemize}
				\item<7-> Integration with Bib\TeX{} for automatic citations
				\item<8-> Efficient workflow for large, structured presentations
			\end{itemize}
		\end{block}
	\end{frame}
	
	\section{Highlighting Texts}
	\begin{frame}{Blocks and Highlights Example}
		Some important text can be \alert{highlighted} to attract attention.
		\begin{block}{Remark}
			This is a normal block. It wraps the text in a box with the same style as
			the presentation. The text after the block title is the content.
		\end{block}
		\begin{alertblock}{Important Theorem}
			This is an alert block. It contrasts with the regular presentation style,
			useful to emphasize key results or theorems.
		\end{alertblock}
		\begin{exampleblock}{Examples}
			This is an example block. It is similar to a normal block, but has a
			different style (less contrasting than alertblock). Perfect for examples.
		\end{exampleblock}
	\end{frame}
	
	\section{Mathematical Concepts}
	\begin{frame}{Definition}
		\begin{definition}
			A \alert{prime number} is a natural number greater than 1 that has no positive divisors other than 1 and itself.
		\end{definition}
	\end{frame}
	
	\begin{frame}{Theorem \& Lemma}
		\begin{theorem}
			\alert{Fundamental Theorem of Arithmetic:} Every integer greater than 1 is either a prime or can be written as a unique product of prime numbers.
		\end{theorem}
		\begin{lemma}
			If a prime $p$ divides a product $ab$, then $p$ divides $a$ or $p$ divides $b$.
		\end{lemma}
	\end{frame}
	\newcounter{algorithmCounter} 
	\setcounter{algorithmCounter}{0}
	\section{Linear Search Algorithm}
	\begin{frame}{Linear Search Algorithm}
		\stepcounter{algorithmCounter}
		\textbf{Input:} Array $A$ of size $n$, key to search \\
		\textbf{Output:} Index of key in $A$, or -1 if not found
		
		\medskip
		for $i \gets 0$ to $n-1$ do\\
		\quad if $A[i] = \text{key}$ then\\
		\quad\quad return $i$;\\
		\quad end if\\
		end for\\
		return -1;
		
		\medskip
		\centering
		\textbf{Algorithm \thealgorithmCounter: Linear Search}
	\end{frame}
	
	\begin{frame}
		\frametitle{Corollary \& Proof}
		\begin{block}{Corollary}
			\emph{The factorization of any integer into primes is unique up to the order of the factors.}
		\end{block}
		
		\begin{block}{Proof}
			Suppose $n$ has two different prime factorizations. By the lemma, every prime in one factorization divides one of the primes in the other factorization. This leads to a contradiction, proving uniqueness.
		\end{block}
	\end{frame}
	\section{Images}
	\begin{frame}{Including an Image in a Block}
		\begin{block}{Example Image}
			\centering
			\includegraphics[scale=0.5]{"/home/bodike/Desktop/B24CS013_Bodike Chaithali_10/overleaf-logo.png"}
			\captionof{figure}{This is a sample image included in a Beamer slide.}
		\end{block}
	\end{frame}
	
	\section{Important Equations}
	\begin{frame}
		\frametitle{Key Equations in Signal Processing}
		\begin{block}{Equation 1: Fourier Transform}
			\[
			X(f) = \int_{-\infty}^{\infty} x(t)e^{-j2\pi ft}dt
			\]
		\end{block}
		\begin{block}{Equation 2: Convolution Integral}
			\[
			y(t) = \int_{-\infty}^{\infty} x(\tau)h(t-\tau)d\tau
			\]
		\end{block}
		\begin{block}{Equation 3: Transfer Function of LTI System}
			\[
			H(f) = \frac{Y(f)}{X(f)} = \int_{-\infty}^{\infty}h(t)e^{-j2\pi ft}dt
			\]
		\end{block}
	\end{frame}
	
	\section{Conclusion}
	\begin{frame}{Conclusion}
		\begin{itemize}[<+->]
			\item \LaTeX{} Beamer is a powerful tool for creating professional presentations.
			\item It supports complex mathematical equations, figures, tables, and
			references.
			\item Overlays and incremental reveals help guide the audience’s focus.
			\item Highly customizable themes, colors, and backgrounds make slides
			visually appealing.
			\item Ideal for academic and technical presentations due to its precision and
			portability.
		\end{itemize}
	\end{frame}
	
	\section{References}
	\begin{frame}[allowframebreaks]
		\frametitle{References}
		This presentation is based on the following references:
		\begin{itemize}
			\item \cite{goossens1994latexcompanion} for basic \LaTeX{} concepts.
			\item \cite{knuth1984tex} for fundamental \TeX{} understanding.
			\item \cite{lamport1994latex} for advanced \LaTeX usage.
			\item \cite{mittelbach2004}for best practices in document preparation.
			\item \cite{kopka2003} for graphics and presentation guidelines in \LaTeX.
			
			\bibliography{ref}
		\end{itemize}
	\end{frame}
	
	\section{Thank You}
	\begin{frame}{Thank You!}
		\centering
		\Huge Thank You!\\[1em] \LARGE Any Questions?\\[1em]
		\includegraphics[scale=0.1]{"/home/bodike/Desktop/B24CS013_Bodike Chaithali_10/iit-bhilai.jpg"}
	\end{frame}
	
\end{document} 